% Auto-generated by Exercise3.jl — do not edit by hand.
\begin{table}[htbp]
  \centering
  \caption{Structural Change and Aggregate Labor Productivity}
  \label{tab:ex3_aggregate}
  \begin{threeparttable}
  \begin{tabular}{l S[table-format=1.2] S[table-format=3.1]}
    \toprule
     & {Avg. annual growth (\%)} & {Productivity index, 2023} \\
     & {1950--2023} & {1950 = 100} \\
    \midrule
    Baseline & 2.04 & 442.9 \\
    Fixed-shares benchmark & 2.13 & 475.1 \\
    Full shutdown & 2.24 & 514.6 \\
    \bottomrule
  \end{tabular}
  \begin{tablenotes}[flushleft]
    \footnotesize
    \RaggedRight
    \item[] \textit{Notes:} Aggregate labor productivity is real value added per production worker.
  Average annual growth is the annualized log change over 1950--2023. Productivity indices are
  normalized to 100 in 1950. The fixed-shares benchmark weights baseline sectoral labor
  productivity growth by 1950 value added shares; its endpoint index is approximated from the
  reported cumulative log growth. The full shutdown sets $\eta_j = 1$ and $\sigma = 1$, with the
  equilibrium re-solved.
  \end{tablenotes}
  \end{threeparttable}
\end{table}
